\section{Results}\label{sec:results}

The results are reported in the order the questions were posed.
Section~\ref{sec:res-grid} gives the grid over the prompt derived corpus, which establishes
what the pipeline achieves when it is trained and tested in the usual way.
Section~\ref{sec:res-denoise} tests the preprocessing step the work being extended relies
on. Section~\ref{sec:res-external} applies the retained weights to human annotated audio
without retraining, which is the test the earlier evaluation did not make.
Section~\ref{sec:res-regime} reports what can and cannot be established from the
classifier derived corpus. Section~\ref{sec:res-limits} states what the results do not
support. Every figure quoted below is computed from the saved per recording decisions by
the scripts released with the code, and the tables are generated from those same files.

\subsection{Detection within the prompt derived corpus}\label{sec:res-grid}

Table~\ref{tab:grid} gives the grid. ResNet18 reaches $99.09 \pm 0.13$ per cent accuracy
on the linear condition, with sensitivity $97.72 \pm 0.62$ and specificity
$99.48 \pm 0.26$, for a balanced accuracy of $98.60 \pm 0.22$ and a correlation
coefficient of $0.974 \pm 0.004$. These are the numbers the literature reports for this
task, obtained here under a partition in which no participant appears on both sides of a
split, so the simplest explanation for a high score is excluded: the networks are not
recognising speakers they were trained on.

\begin{table}[htbp]
\caption{Detection on the prompt derived corpus, over held out speakers. Each row is one architecture under one input condition, averaged over the seeds in the third column, with the standard deviation across those seeds. Sensitivity is the proportion of coughs recovered and specificity the proportion of non-coughs rejected; both are reported because the class prior makes accuracy alone uninformative.}\label{tab:grid}
\small
\begin{tabular}{@{}llrcccc@{}}
\toprule
Architecture & Input & Seeds & Accuracy & Sensitivity & Specificity & MCC \\
\midrule
ResNet18 & linear & 5 & $99.09 \pm 0.13$ & $97.72 \pm 0.62$ & $99.48 \pm 0.26$ & $0.974 \pm 0.004$ \\
 & wav+linear & 5 & $98.80 \pm 0.24$ & $97.14 \pm 0.53$ & $99.27 \pm 0.30$ & $0.965 \pm 0.007$ \\
 & mel & 5 & $98.71 \pm 0.20$ & $97.19 \pm 0.81$ & $99.14 \pm 0.10$ & $0.963 \pm 0.006$ \\
 & wav+mel & 5 & $98.65 \pm 0.49$ & $97.37 \pm 0.53$ & $99.02 \pm 0.74$ & $0.961 \pm 0.014$ \\
\addlinespace
MobileNetV2 & mel & 5 & $98.86 \pm 0.33$ & $98.35 \pm 0.34$ & $99.00 \pm 0.49$ & $0.967 \pm 0.009$ \\
 & wav+mel & 4 & $98.96 \pm 0.24$ & $97.88 \pm 0.39$ & $99.27 \pm 0.32$ & $0.970 \pm 0.007$ \\
\addlinespace
Compact & linear & 5 & $96.43 \pm 0.38$ & $93.35 \pm 1.60$ & $97.31 \pm 0.74$ & $0.898 \pm 0.010$ \\
 & mel & 5 & $96.46 \pm 0.34$ & $92.46 \pm 2.24$ & $97.60 \pm 0.83$ & $0.898 \pm 0.009$ \\
 & wav+mel & 5 & $92.57 \pm 8.29$ & $73.93 \pm 41.40$ & $97.89 \pm 1.55$ & $0.715 \pm 0.400$ \\
\botrule
\end{tabular}
\end{table}



The compact network reaches $96.43 \pm 0.38$ per cent on the same condition, and the gap
between it and ResNet18 is larger in sensitivity, $93.35 \pm 1.60$ against
$97.72 \pm 0.62$, than in specificity, where the two are within a point and a half of each
other. Capacity is therefore being spent on recovering coughs rather than on rejecting the
negatives, which is consistent with a negative class that is not difficult.

One result in the table is a failure rather than a measurement, and it is reported because
it bears on the stability of the smaller network. The compact network on the wavelet
shrunk mel condition averages $73.93 \pm 41.40$ per cent sensitivity. A standard deviation
of that size on a quantity bounded by one hundred is not a spread about a mean; it records
that the seeds did not all produce the same kind of model. One seed converged to the
majority class and stayed there, which on this prior produces an accuracy equal to the
non-cough frequency and a sensitivity of zero. The trainer reports any optimiser step
abandoned for a non-finite gradient, and none were abandoned in that run, so this is an
optimisation outcome rather than the arithmetic failure described in
Section~\ref{sec:setup-train}. The relevant observation is that it happened under the
condition combining the mel axis with shrinkage and under no other, and only to the
network with the least capacity.

DenseNet201 was run on the two mel conditions at three seeds, because its memory
requirement makes it several times more expensive per run than the others, and what it
shows is worth that cost. On the mel axis it reaches $98.66 \pm 0.31$ per cent balanced
accuracy, against MobileNetV2's $98.68 \pm 0.17$ and ResNet18's $98.17 \pm 0.41$. The
deepest network in the comparison is therefore indistinguishable from the smallest
pretrained one, and its margin over the middle one is half a point, which is inside the
bound Section~\ref{sec:setup-metrics} sets for an effect that may be claimed at all. The
separation that does appear is the three to four points between each of the three
pretrained networks and the one trained from random initialisation on the same axis. On
this task, at this label quality, it is pretrained initialisation and not capacity that
distinguishes one architecture from another.

\subsection{Whether wavelet shrinkage helps}\label{sec:res-denoise}

Table~\ref{tab:paired} tests the shrinkage on paired decisions. The result is that it does
not help, and that where the effect reaches significance it is negative. \figref{fig:effect} sets each effect against the spread across seeds it would have to clear to be claimable from single runs, and most of them fall on the wrong side of that line.

\begin{table*}[tbp]
\caption{Wavelet shrinkage tested on paired decisions over the prompt derived corpus.}\label{tab:paired}
\small
\begin{tabular}{@{}llrrrllc@{}}
\toprule
Architecture & Comparison & Seeds & $b$ & $c$ & $p$ & $p$ adj.\ & Accuracy change \\
\midrule
ResNet18 & linear $\rightarrow$ wav+linear & 5 & 66 & 37 & $0.0055$ & $0.0277$ & $-0.29 \pm 0.19$ \\
ResNet18 & mel $\rightarrow$ wav+mel & 5 & 71 & 65 & $0.668$ & $1.000$ & $-0.06 \pm 0.39$ \\
MobileNetV2 & linear $\rightarrow$ wav+linear & 5 & 71 & 50 & $0.0686$ & $0.206$ & $-0.21 \pm 0.24$ \\
MobileNetV2 & mel $\rightarrow$ wav+mel & 5 & 53 & 58 & $0.704$ & $1.000$ & $+0.05 \pm 0.40$ \\
DenseNet201 & mel $\rightarrow$ wav+mel & 3 & 50 & 29 & $0.0238$ & $0.0953$ & $-0.35 \pm 0.52$ \\
Compact & linear $\rightarrow$ wav+linear & 5 & 196 & 115 & $<10^{-4}$ & $<10^{-4}$ & $-0.80 \pm 0.65$ \\
Compact & mel $\rightarrow$ wav+mel & 5 & 525 & 133 & $<10^{-4}$ & $<10^{-4}$ & $-3.89 \pm 8.25$ \\
\bottomrule
\end{tabular}
\end{table*}



\begin{figure}[tbp]
\centering
\includegraphics[width=0.98\columnwidth]{effect.pdf}
\caption{Each shrinkage effect against the seed spread it would have to clear.}\label{fig:effect}
\end{figure}

For ResNet18 on the linear axis, $66$ recordings are classified correctly without
shrinkage and incorrectly with it, against $37$ in the reverse direction. McNemar's test
on those discordant pairs gives $p = 0.0055$, which survives the step down correction over
the family in Table~\ref{tab:paired}, and the accuracy change across seeds is
$-0.29 \pm 0.19$ points. The effect is therefore real, small, and in the direction
opposite to the one claimed for the method. On the mel axis the same comparison is
indistinguishable from no effect, $71$ against $65$ discordant recordings and
$p = 0.668$ before any correction. For the compact network on the mel axis the test
rejects decisively, but the accuracy change is $-3.89 \pm 8.25$ points, and a spread four
times the size of the effect is the degenerate seed discussed above rather than a property
of shrinkage. The same network on the linear axis carries the same conclusion without that
complication: $196$ recordings against $115$, a decisive rejection, and a change of
$-0.80 \pm 0.65$ points with no degenerate seed behind it.

DenseNet201 contributes the one comparison that shows what the correction is for. Its
single test gives $p = 0.0238$, which any convention that ignores the family a test
belongs to would report as an effect, and $0.0953$ after the step down, which would not.
Had that architecture been run on its own, the same decisions would have been published as
a significant result about shrinkage on the mel axis, and over the family they do not
support one.

Section~\ref{sec:design-denoise} supplies the mechanism, and it is worth restating here
because it determines how far these results generalise. The universal threshold is set
from an estimate of the noise level in the finest detail coefficients. On recordings
already as clean as these it falls far below the signal, and the measured consequence is
that a median of $0.67$ per cent of signal power is removed, with quartiles at $0.08$ and
$2.28$ per cent. The two input conditions are therefore very nearly the same images, and
the null result is what should be expected rather than a surprise. This is a statement
about applying this family of shrinkage rules to corpora of this quality; it is not
evidence that denoising cannot help audio that is genuinely noisy.

\subsection{Transfer to human annotated audio}\label{sec:res-external}

The weights retained from the grid were applied, without retraining and without adaptation
of any kind, to the exhaustively verified evaluation split of the human annotated corpus
at the same class prior. Table~\ref{tab:external} gives the result, its rows corresponding one to one with
those of Table~\ref{tab:grid}, and \figref{fig:transfer} plots it against the in corpus
figure for the same weights. In that figure the dotted line marks equal performance and
the error bars are one standard deviation across seeds on each axis; the input
conditions run for each architecture are plotted without being labelled separately,
because they fall on top of one another, the choice of frequency axis and the presence
or absence of shrinkage moving the result by far less than the change of corpus does.
The fall itself is asymmetric: specificity is substantially retained while sensitivity
is not, so the detectors go on rejecting non-coughs and stop recovering coughs.

\begin{table}[htbp]
\caption{The same saved weights applied to human annotated audio at the same class prior, with no retraining and no adaptation of any kind. The rows correspond one to one with those of Table~\ref{tab:grid}. The fall is asymmetric: specificity is substantially retained while sensitivity is not, so the detectors continue to reject non-coughs and stop recovering coughs.}\label{tab:external}
\small
\begin{tabular}{@{}llrcccc@{}}
\toprule
Architecture & Input & Seeds & Accuracy & Sensitivity & Specificity & MCC \\
\midrule
ResNet18 & linear & 5 & $78.03 \pm 1.62$ & $56.04 \pm 4.04$ & $84.32 \pm 3.24$ & $0.391 \pm 0.012$ \\
 & wav+linear & 5 & $79.33 \pm 1.19$ & $53.21 \pm 1.27$ & $86.81 \pm 1.65$ & $0.402 \pm 0.022$ \\
 & mel & 5 & $78.70 \pm 2.31$ & $57.55 \pm 3.89$ & $84.76 \pm 3.88$ & $0.411 \pm 0.031$ \\
 & wav+mel & 5 & $77.56 \pm 3.18$ & $59.62 \pm 2.78$ & $82.70 \pm 4.34$ & $0.401 \pm 0.056$ \\
\addlinespace
MobileNetV2 & linear & 2 & $76.26 \pm 2.08$ & $62.74 \pm 2.00$ & $80.14 \pm 3.25$ & $0.392 \pm 0.023$ \\
 & mel & 5 & $74.92 \pm 3.98$ & $66.04 \pm 4.11$ & $77.46 \pm 5.97$ & $0.392 \pm 0.054$ \\
 & wav+mel & 5 & $77.65 \pm 3.00$ & $57.92 \pm 4.30$ & $83.30 \pm 4.89$ & $0.396 \pm 0.038$ \\
\addlinespace
Compact & linear & 5 & $68.11 \pm 1.55$ & $45.66 \pm 4.88$ & $74.54 \pm 3.26$ & $0.184 \pm 0.015$ \\
 & mel & 5 & $73.36 \pm 5.67$ & $38.49 \pm 9.70$ & $83.35 \pm 10.00$ & $0.235 \pm 0.049$ \\
 & wav+mel & 5 & $70.34 \pm 8.57$ & $36.79 \pm 21.85$ & $79.95 \pm 16.25$ & $0.161 \pm 0.106$ \\
\bottomrule
\end{tabular}
\end{table}



ResNet18 on the linear condition falls from $98.60$ per cent balanced accuracy to
$70.18 \pm 0.44$. The fall is not uniform across the two classes, and the asymmetry is the
substantive finding. Specificity declines from $99.48 \pm 0.26$ to $84.32 \pm 3.24$, which
is a real loss but leaves the detector still rejecting most non-coughs. Sensitivity falls
from $97.72 \pm 0.62$ to $56.04 \pm 4.04$. A detector that recovered better than ninety
seven coughs in a hundred within the corpus it was fitted to recovers barely more than
half of them once the recordings, the microphones, the annotators and the definition of
the negative class all change together.

\begin{figure*}[tbp]
\centering
\includegraphics[width=0.74\textwidth]{transfer.pdf}
\caption{In corpus against out of corpus balanced accuracy for the same saved
weights.}\label{fig:transfer}
\end{figure*}

A fall of this size has two possible causes that call for quite different responses, and
they can be told apart. The networks may have stopped ranking coughs above non-coughs, in
which case the representation does not transfer and nothing applied after the fact will
recover it. Or they may still rank them correctly while the decision threshold inherited
from training now sits in the wrong place, in which case the loss is a calibration problem
and a practitioner holding a little labelled data from the target domain can fix it
cheaply. The area under the receiver operating characteristic distinguishes the two,
because it depends on the ordering of the scores alone and not on where the threshold
falls. Table~\ref{tab:threshold} reports it.

\begin{table*}[tbp]
\caption{Discrimination against calibration on the external set.}\label{tab:threshold}
\small
\begin{tabular}{@{}llcccc@{}}
\toprule
& & \multicolumn{2}{c}{Area under the curve} & \multicolumn{2}{c}{External balanced accuracy} \\
\cmidrule(lr){3-4}\cmidrule(lr){5-6}
Architecture & Input & In corpus & External & At $0.5$ & Best possible \\
\midrule
ResNet18 & linear & $0.998 \pm 0.002$ & $0.749 \pm 0.018$ & $70.18 \pm 0.44$ & $71.90 \pm 0.91$ \\
 & wav+linear & $0.998 \pm 0.001$ & $0.730 \pm 0.007$ & $70.01 \pm 0.81$ & $71.13 \pm 0.52$ \\
 & mel & $0.999 \pm 0.000$ & $0.768 \pm 0.016$ & $71.15 \pm 1.01$ & $72.34 \pm 1.68$ \\
 & wav+mel & $0.998 \pm 0.001$ & $0.754 \pm 0.024$ & $71.16 \pm 2.08$ & $72.31 \pm 1.35$ \\
\addlinespace
MobileNetV2 & linear & $0.998 \pm 0.000$ & $0.768 \pm 0.007$ & $71.88 \pm 1.25$ & $73.30 \pm 0.73$ \\
 & wav+linear & $0.998 \pm 0.001$ & $0.739 \pm 0.014$ & $69.05 \pm 1.45$ & $70.61 \pm 1.95$ \\
 & mel & $0.998 \pm 0.000$ & $0.766 \pm 0.025$ & $71.75 \pm 1.93$ & $72.90 \pm 1.62$ \\
 & wav+mel & $0.998 \pm 0.001$ & $0.748 \pm 0.023$ & $70.61 \pm 1.17$ & $71.64 \pm 1.65$ \\
\addlinespace
DenseNet201 & mel & $0.998 \pm 0.001$ & $0.767 \pm 0.022$ & $71.57 \pm 0.96$ & $72.96 \pm 0.93$ \\
 & wav+mel & $0.998 \pm 0.002$ & $0.736 \pm 0.021$ & $69.27 \pm 4.09$ & $70.99 \pm 3.07$ \\
\addlinespace
Compact & linear & $0.992 \pm 0.001$ & $0.648 \pm 0.012$ & $60.10 \pm 1.06$ & $62.23 \pm 1.59$ \\
 & wav+linear & $0.988 \pm 0.002$ & $0.650 \pm 0.019$ & $60.45 \pm 1.87$ & $62.81 \pm 1.52$ \\
 & mel & $0.993 \pm 0.000$ & $0.687 \pm 0.017$ & $60.92 \pm 1.01$ & $64.36 \pm 1.58$ \\
 & wav+mel & $0.985 \pm 0.012$ & $0.651 \pm 0.028$ & $58.37 \pm 5.18$ & $62.10 \pm 2.02$ \\
\bottomrule
\end{tabular}
\end{table*}



The answer is the unwelcome one. For ResNet18 on the linear condition the area under the
curve falls from $0.998$ within the corpus to $0.749$ outside it. The ordering itself has
degraded, and the threshold is very nearly where it ought to be: balanced accuracy at the
inherited threshold is $70.18$ per cent, while the best value attainable at any threshold
on this set is $71.90$. Recalibration could therefore recover about $1.7$ points, and even
that figure is optimistic, because the threshold achieving it was chosen on the same
recordings it is scored on and no procedure could have selected it in advance. The same
pattern holds for every architecture and condition in the table. The obvious and cheapest
remedy for a domain shift, collecting a little target data and moving the operating point,
would recover almost none of what was lost here.

Two further observations follow from the external table. The first is that the only
separation between architectures that was worth anything inside the corpus is also the
only one that survives leaving it, and it does not lie between the pretrained networks.
ResNet18, MobileNetV2 and DenseNet201 reach $71.15 \pm 1.01$, $71.75 \pm 1.93$ and
$71.57 \pm 0.96$ per cent balanced accuracy on the mel axis outside the corpus, which
places all three inside one another's spread, as they were inside it. The compact network,
roughly three points behind ResNet18 in corpus, is roughly ten points behind it out of
corpus, at $60.10 \pm 1.06$ on the linear condition. Whatever pretraining supplied that
the smaller network had to find in the data alone, part of it transfers, and capacity
beyond that supplies nothing in either regime. The second is that the input conditions remain
indistinguishable from one another under the change of corpus. Shrinkage gains
$1.30 \pm 1.36$ accuracy points on the linear axis externally, which is a change smaller
than its own spread across seeds, so the external set does not rescue the preprocessing
step either.

The degenerate seed reappears here and is worth following, because it shows what a single
run can conceal. The compact network on the wavelet shrunk mel condition records, for one
seed, an accuracy of $77.73$ per cent against a balanced accuracy of exactly $50.00$. The
first number is simply the proportion of non-coughs in the external set. The model
assigned every one of the $476$ recordings to the majority class, and an evaluation
reporting accuracy alone would have recorded it as the second best run of that condition.

Because the external negatives are annotated against the ontology, it is possible to ask
not merely how many false positives there are but what produces them.
Table~\ref{tab:confusion} gives the rate at which clips carrying each annotated class are
assigned to the cough class, pooled over every ResNet18 run, and \figref{fig:falsepos} orders the classes by that rate. A clip may carry several labels, so the rows overlap and their counts do not sum to the size of the negative class.

\begin{table}[htbp]
\caption{What the detectors call a cough, by annotated class, pooled over every ResNet18 run on the external set.}\label{tab:confusion}
\small
\begin{tabular}{@{}lrrc@{}}
\toprule
Annotated class & Decisions & Called cough & Per cent \\
\midrule
Sneeze & 360 & 244 & $67.8$ \\
Burping & 440 & 149 & $33.9$ \\
Whispering & 520 & 139 & $26.7$ \\
Gasp & 200 & 41 & $20.5$ \\
Sigh & 280 & 56 & $20.0$ \\
Breathing & 800 & 132 & $16.5$ \\
Yell & 280 & 40 & $14.3$ \\
Male speech & 2320 & 316 & $13.6$ \\
Speech & 4180 & 492 & $11.8$ \\
Laughter & 1400 & 161 & $11.5$ \\
Female speech & 1460 & 143 & $9.8$ \\
Child speech & 440 & 34 & $7.7$ \\
Chewing & 360 & 22 & $6.1$ \\
\addlinespace
Cough, as recall & 2120 & 1200 & $56.6$ \\
\bottomrule
\end{tabular}
\end{table}



\begin{figure}[tbp]
\centering
\includegraphics[width=0.98\columnwidth]{falsepos.pdf}
\caption{The rate at which each annotated class is assigned to the cough class.}\label{fig:falsepos}
\end{figure}

The errors are not spread evenly over the negative class, and the concentration is
severe. Clips labelled \textsf{Sneeze} are called coughs in $67.8$ per cent of decisions.
That is a higher rate than the detectors achieve on clips that really are coughs, where
recall over the same runs is $56.6$ per cent. On this evidence a sneeze is more likely to
be reported as a cough than a cough is. Burping follows at $33.9$ per cent, whispering at
$26.7$, gasping at $20.5$ and sighing at $20.0$, while breathing sits at $16.5$ and the
speech classes fall between $7.7$ and $13.6$.

The ordering is interpretable, and it is the one the construction of the prompted corpus
predicts. The sounds most readily confused with coughing are the other forced
expirations. A sneeze, a burp, a gasp and a sigh are each a sudden release of air through
the vocal tract, distinguished from a cough by details a detector has no reason to encode
if it was never required to. The negative class of the prompted corpus is breathing,
sustained vowels and counting, and not one of those is a forced expiration. A model
trained against it is never asked to separate coughing from sneezing, and the external
evaluation finds that it cannot. This is the mechanism behind the loss of discrimination
reported above, and it indicates what a corpus intended to support cough detection would
need to contain: not more recordings, but negatives that are hard in the particular way
that matters.

\subsection{What the classifier derived corpus can establish}\label{sec:res-regime}

The third regime is reported separately because the quantity it measures is different in
kind. A model trained and tested on labels thresholded from a detector score is measured on
its agreement with that detector. A high figure there is evidence that two classifiers
agree, which is not evidence that a cough is present, and no amount of care in the
partitioning changes what is being measured.

At the matched prior this corpus supplies $8398$ recordings, of which $1880$ carry a cough
label, and its test partition holds $1259$ of them. Within that regime a ResNet18 reaches
$98.41 \pm 0.24$ per cent balanced accuracy on linear spectrograms, against $98.60 \pm
0.22$ on the prompt derived corpus under the same architecture and the same input
condition. Table~\ref{tab:regime} gives the four conditions. The resemblance between those
two figures is the most instructive thing about them: a reader given only the headline
number could not tell which of the two experiments produced it, although one measures
recovery of a prompted behaviour from held out speakers and the other measures agreement
with a detector. Nothing in the quantity itself marks the difference, which is why the
regime is named in the caption rather than left to be inferred from the corpus name.

\begin{table*}[tbp]
\caption{Detection on the classifier derived corpus at a matched class prior. What is measured here is agreement with the detector that produced the labels, not whether a cough is present, and the figures should be read as a property of that regime.}\label{tab:regime}
\small
\begin{tabular}{@{}llrcccc@{}}
\toprule
Architecture & Input & Seeds & Accuracy & Sensitivity & Specificity & MCC \\
\midrule
ResNet18 & linear & 5 & $98.94 \pm 0.14$ & $97.45 \pm 0.68$ & $99.37 \pm 0.29$ & $0.969 \pm 0.004$ \\
 & wav+linear & 5 & $98.25 \pm 0.51$ & $97.87 \pm 1.00$ & $98.36 \pm 0.85$ & $0.951 \pm 0.013$ \\
 & mel & 5 & $98.57 \pm 0.32$ & $98.09 \pm 1.05$ & $98.71 \pm 0.52$ & $0.960 \pm 0.009$ \\
 & wav+mel & 5 & $98.60 \pm 0.25$ & $97.38 \pm 0.48$ & $98.96 \pm 0.36$ & $0.960 \pm 0.007$ \\
\bottomrule
\end{tabular}
\end{table*}



The paired test behaves on this corpus much as it does on the other. Shrinkage applied
before a linear spectrogram costs $0.68 \pm 0.38$ accuracy points, with $79$ recordings
lost against $36$ gained and an adjusted $p$ of $0.0002$; applied before a mel spectrogram
it changes nothing measurable, $51$ against $53$ and $p = 0.922$. The direction agrees
with the prompt derived corpus in both pairs and the magnitude is larger in the first,
which is consistent with shrinkage removing detail that the linear axis still resolves and
the mel axis has already averaged together. What a loss of agreement means here needs
stating, though. These labels were assigned to the recordings as distributed rather than to
shrunk versions of them, so a processing step that moves a model away from those labels has
not thereby been shown to move it away from the coughs.

This corpus also cannot support the partitioning the other two receive. It distributes no
participant identifier, so recordings contributed by one person cannot be kept to one side
of a split, and the figures in Table~\ref{tab:regime} are accordingly not protected against
a contributor appearing in both training and test. That is a second reason, independent of
the labels, to read them as a property of the regime rather than as a detection result.

Two properties of the corpus bound what can be done with it, and both are settled from the
published metadata without reference to any model. The first concerns reproducing the
selection used by the work being extended, which reports $15{,}619$ cough and $7{,}261$
non-cough recordings. The lower count is reachable: a threshold of $0.1271$ on the detector
score yields exactly $7{,}261$. The upper count is reachable by no threshold on that
column. Six recordings carry the score $0.8961$ exactly, so the number at or above a
threshold in that neighbourhood steps from $15{,}617$ directly to $15{,}623$, and
$15{,}619$ lies strictly between two consecutive attainable values. Some further criterion
was applied that the original does not record. This work therefore uses round thresholds
and states them rather than tuning towards a near miss.

The second concerns the expert annotations. The corpus carries judgements from clinical
annotators on a subset of its recordings, and it is tempting to use agreement on that
subset to estimate the label error rate of the corpus as a whole. That inference requires
the annotated subset to be a random sample, and it is not. Of $34{,}434$ recordings,
$2{,}890$ carry a judgement from at least one of the four annotators, which is $8.4$ per
cent of the corpus. Their detector scores average $0.8954$ against $0.6179$ for the
remainder, and $85.8$ per cent of them score at or above $0.8$ against $50.9$ per cent of
the unannotated recordings. The recordings sent for
annotation were drawn overwhelmingly from the region where the detector was already
confident. Agreement measured there cannot be extrapolated to the recordings nearest the
threshold, which are precisely the ones whose labels are least secure, so the label error
rate of this corpus is not estimable from what it publishes.

\subsection{What these results do not establish}\label{sec:res-limits}

The external evaluation rests on a single corpus, and on the $476$ recordings its
evaluation split supplies at the matched prior, of which $106$ are coughs. The direction
and the approximate size of the fall are secure, since the interval estimates do not come
close to the in corpus figures, but a set of this size cannot pin the magnitude down
finely, and one external corpus cannot distinguish a general failure to transfer from a
particular mismatch between these two collections. A second human annotated corpus would
separate those, and none was available that met the conditions used here.

The three regimes differ in more than the provenance of their labels, and the matching
performed here does not remove every difference. The class prior is matched by
construction and the preprocessing, architecture and protocol are held fixed, but the
recording conditions, the devices and the populations differ between corpora and cannot be
equalised after the fact. What the comparison isolates is the combined effect of label
provenance and negative class composition; it does not separate those two from each other.
Doing so would require one corpus annotated twice, once by prompt and once by a listener,
over the same recordings, which does not currently exist.

The partition over the classifier derived corpus is not speaker disjoint, because that
corpus publishes no participant identifier. Any contributor who submitted more than once
may appear on both sides of a split, and the figures for that regime should be read with
that in mind. The same limitation does not affect the prompt derived corpus, where the
partition is built over participants, nor the human annotated split, which is never
trained on.

Finally, the architectures examined here are standard image classifiers applied to
spectrograms, trained for at most eight epochs with a fixed learning rate and no
augmentation. A longer schedule, a search over learning rates, or augmentation designed for
audio might all raise the in corpus figures. None of them is likely to change the finding
that matters, because the gap reported here is not between one training recipe and another
but between performance measured inside a corpus and performance measured outside it, and
a recipe that improves the former has no particular reason to improve the latter.
